\documentclass[a4paper,12pt]{article}
\usepackage[utf8]{inputenc}
\usepackage[spanish]{babel}
\usepackage{geometry}
\usepackage{xcolor}
\usepackage{listings}
\usepackage{hyperref}

% Configuración de márgenes
\geometry{top=2.5cm, bottom=2.5cm, left=2.5cm, right=2.5cm}

% Configuración de colores para código
\definecolor{codegreen}{rgb}{0,0.6,0}
\definecolor{codegray}{rgb}{0.5,0.5,0.5}
\definecolor{codepurple}{rgb}{0.58,0,0.82}
\definecolor{backcolour}{rgb}{0.95,0.95,0.92}

% Define YAML language for listings
\lstdefinelanguage{yaml}{
  keywords={services, ports, depends_on, build, app, db, payments},
  keywordstyle=\color{magenta},
  ndkeywords={true, false, null},
  ndkeywordstyle=\color{blue},
  identifierstyle=\color{black},
  sensitive=false,
  comment=[l]{\#},
  commentstyle=\color{codegreen},
  stringstyle=\color{codepurple},
  morestring=[b]',
  morestring=[b]"
}

% Estilo para código Go
\lstdefinestyle{mystyle}{
    backgroundcolor=\color{backcolour},   
    commentstyle=\color{codegreen},
    keywordstyle=\color{magenta},
    numberstyle=\tiny\color{codegray},
    stringstyle=\color{codepurple},
    basicstyle=\ttfamily\footnotesize,
    breakatwhitespace=false,         
    breaklines=true,                 
    captionpos=b,                    
    keepspaces=true,                 
    numbers=left,                    
    numbersep=5pt,                  
    showspaces=false,                
    showstringspaces=false,
    showtabs=false,                  
    tabsize=2
}

\lstset{style=mystyle}

\title{\textbf{Compendio de Desarrollo Backend}\\ \large Go, PostgreSQL y Docker}
\author{Inventory Manager Project}
\date{\today}

\begin{document}

\maketitle
\tableofcontents
\newpage

\section{Arquitectura y Estructura del Proyecto}

Pasamos de un script monolítico a una \textbf{Arquitectura Modular}. Esto separa responsabilidades, haciendo el código limpio y mantenible.

\begin{itemize}
    \item \textbf{\texttt{main.go} (El Director):} Solo se encarga de recibir peticiones HTTP (rutas) y decidir qué función llamar. No contiene lógica SQL.
    \item \textbf{\texttt{database/} (El Almacén):} Paquete separado donde reside la conexión y lógica SQL.
    \item \textbf{Visibilidad (Scope):}
    \begin{itemize}
        \item \textbf{Mayúscula} (ej. \texttt{InitDB}, \texttt{Producto}): Público/Exportado.
        \item \textbf{Minúscula} (ej. \texttt{crearTablas}): Privado/Interno.
    \end{itemize}
\end{itemize}

\section{Infraestructura: Docker y Docker Compose}

Todo corre aislado en contenedores, sin instalar dependencias en el sistema operativo host.

\begin{itemize}
    \item \textbf{Contenedores:} Entornos ligeros para ejecutar la aplicación y la base de datos.
    \item \textbf{\texttt{docker-compose.yml}:} Orquestador que define los servicios:
    \begin{enumerate}
        \item \textbf{\texttt{app}:} Código Go (Puerto 8080).
        \item \textbf{\texttt{db}:} PostgreSQL.
    \end{enumerate}
    \item \textbf{Red Interna:} Los contenedores se ven por nombre de servicio (host: \texttt{"db"}).
    \item \textbf{Volúmenes (Persistencia):}
    \begin{verbatim}
volumes:
  - postgres_data:/var/lib/postgresql/data
    \end{verbatim}
    Vital para que los datos sobrevivan al reiniciar los contenedores.
\end{itemize}

\section{El Lenguaje: Go (Golang)}

\subsection{Estructuras (Structs) y Tags}
El plano de los datos. Los "Tags" definen cómo se visualizan en JSON.

\begin{lstlisting}[language=Go]
type Producto struct {
    ID     int     `json:"id"`
    Nombre string  `json:"nombre"`
    Stock  int     `json:"stock"`
}
\end{lstlisting}

\subsection{Manejo de Errores}
Go no usa \texttt{try/catch}. Se verifica el error explícitamente.

\begin{lstlisting}[language=Go]
resultado, err := funcionPeligrosa()
if err != nil {
    log.Fatal(err) // Si falla, detenemos la ejecucion
}
\end{lstlisting}

\subsection{Identificador en Blanco (\_)}
Se usa cuando una función devuelve un valor que no necesitamos.
\begin{lstlisting}[language=Go]
// Solo importa el error, ignoramos el resultado sql.Result
_, err := DB.Exec(query) 
\end{lstlisting}

\section{Base de Datos: SQL y PostgreSQL}

\subsection{Tipos de Consultas}
\begin{itemize}
    \item \textbf{\texttt{DB.Exec}:} Para INSERT, UPDATE, DELETE, CREATE.
    \item \textbf{\texttt{DB.Query}:} Para SELECT que devuelve múltiples filas.
    \item \textbf{\texttt{DB.QueryRow}:} Para SELECT de un solo registro (por ID).
\end{itemize}

\subsection{Seguridad y Placeholders}
Uso de \texttt{\$1, \$2} para evitar Inyección SQL.

\begin{lstlisting}[language=Go]
query := "SELECT * FROM productos WHERE id = $1"
DB.QueryRow(query, idVariable)
\end{lstlisting}

\subsection{Integridad Referencial}
Uso de Foreign Keys para relacionar tablas (ej. Ordenes y Productos).
\begin{lstlisting}[language=SQL]
producto_id INTEGER REFERENCES productos(id)
\end{lstlisting}

\section{Integración: HTTP y API REST}

\begin{itemize}
    \item \textbf{GET:} Leer datos. Parámetros en URL (\texttt{/productos/5}).
    \item \textbf{POST:} Crear/Accionar. Datos en Body (JSON).
\end{itemize}

\subsection{Decodificar (Leer JSON)}
\begin{lstlisting}[language=Go]
var p Producto
json.NewDecoder(r.Body).Decode(&p)
\end{lstlisting}

\subsection{Codificar (Responder JSON)}
\begin{lstlisting}[language=Go]
w.Header().Set("Content-Type", "application/json")
json.NewEncoder(w).Encode(p)
\end{lstlisting}

\section{Lógica de Negocio (Ejemplo Venta)}

Flujo transaccional implementado:
\begin{enumerate}
    \item \textbf{Recibir:} Orden con ID producto y Cantidad.
    \item \textbf{Validar:} Existencia y Stock suficiente.
    \item \textbf{Calcular:} Total = Precio $\times$ Cantidad (Backend autoritativo).
    \item \textbf{Actualizar:} Restar stock (\texttt{UPDATE}).
    \item \textbf{Registrar:} Guardar orden (\texttt{INSERT}).
    \item \textbf{Responder:} Devolver objeto completo.
\end{enumerate}

\section{Cheat Sheet de Comandos}

\begin{description}
    \item[docker compose up --build] Reconstruye y levanta los contenedores. Usar ante cambios de código.
    \item[git checkout -b rama] Crea y cambia a una nueva rama.
    \item[go mod init] Inicia un nuevo proyecto Go.
    \item[defer rows.Close()] Cierra la conexión a la DB al finalizar la función.
\end{description}




\newpage
\section{Arquitectura de Microservicios}

Evolucionamos el sistema monolítico a una arquitectura distribuida, separando la lógica de cobros en un servicio independiente.

\subsection{Comunicación entre Contenedores}
Utilizamos la red interna de Docker. Los servicios se resuelven por su nombre en \texttt{docker-compose.yml}, no por IP ni localhost.

\begin{itemize}
    \item \textbf{Inventario:} \texttt{http://app:8080}
    \item \textbf{Pagos:} \texttt{http://payments:8081}
\end{itemize}

\subsection{Cliente HTTP en Go (Hacer Peticiones)}
El backend ya no solo recibe peticiones (Servidor), ahora también las hace (Cliente).

\begin{lstlisting}[language=Go]
// 1. Empaquetar datos (Struct -> JSON Bytes)
jsonData, _ := json.Marshal(solicitud)

// 2. Crear buffer (Reader)
body := bytes.NewBuffer(jsonData)

// 3. Realizar el POST a la URL interna de Docker
// Notar el host "payments" en lugar de "localhost"
resp, err := http.Post("http://payments:8081/cobrar", "application/json", body)
\end{lstlisting}

\subsection{Docker Compose Multi-Servicio}
Orquestación de múltiples contenedores simultáneos.

\begin{lstlisting}[language=yaml, style=mystyle]
services:
  app:
    ports: ["8080:8080"]
    depends_on: [db, payments] # Esperar a los vecinos
  
  payments:
    build: ./payment-service
    ports: ["8081:8081"]
\end{lstlisting}

\subsection{Flujo de Venta Distribuido}
\begin{enumerate}
    \item Cliente envía orden a \texttt{POST /vender}.
    \item Inventario calcula total (Precio $\times$ Cantidad).
    \item \textbf{Bloqueante:} Inventario llama a Pagos y espera 2 segundos.
    \item Si Pagos responde \texttt{200 OK}, se guarda en DB.
    \item Si Pagos falla, se cancela la transacción (Integridad).
\end{enumerate}





\end{document}